\section{Alcance de la aritmética de intervalos}\label{sec:numcert}
¿Qué anchura puede garantizarse para $\cerr$ al sustituir las estimaciones de malla por cotas
rigurosas? El primer error que se propaga a todas las filas es el del perfil. Resolver sus
ecuaciones en punto flotante proporciona un candidato, y el error del primero dominaba el ensayo
de las filas explícitas. Para el nuevo candidato se usó un polinomio por tramos con coeficientes
racionales, de grado 14 y con 16 paneles por unidad. Se encerró su residuo en las ecuaciones en
todo el dominio. La coercividad en variables logarítmicas y la identificación con la solución
física convierten ese residuo en radios uniformes de $2{,}086\,10^{-15}$ para el factor axial y
$4{,}248\,10^{-16}$ para la vacancia [O]. Cada fila debe usar los radios del candidato que evalúa.

\paragraph{Del perfil a los valores nodales [O].}
Los 43 nodos distintos del cálculo comprenden 26 nodos reales de Gauss y 17 representantes del
círculo de 32 puntos. Se encerraron $G_{\rm cof}, f_{4,\rm expl}, f_{2F}$ y el momento $T_1$ en los 43.
La contribución de cuadrados $C_{3,\rm sq}$ también tiene encierros en estos puntos. En el eje real se
encerraron además momentos anidados y la parte explícita de $\mathcal V$. Para las dos coincidencias
$\Delta_1,\Delta_2$ se obtuvieron cotas de norma. Alcanzar las
tolerancias que usa la cuadratura requiere también controlar esas dos contribuciones y las filas
que contienen la resolvente con suficiente precisión.

\paragraph{Valores de frontera [O].}
Para $S$, la simetría y la división en los pliegues del perfil reducen la integral a 27 tetraedros. Las
envolventes polinómicas, los nodos de Gauss verificados, un resto analítico y el redondeo
exterior~\cite{mkc} dan
\[
 15{,}6334582797746\le S\le 15{,}6334582811044 .
\]
Para $P_T$, los productos intervalares encierran las potencias del núcleo centrado de torres y su
variación de medida. El plegado por reflexión, con 2048 celdas por semieje y potencias hasta el orden
14, da
\[
 0{,}444317034618897\le P_T\le 0{,}446141238271620 .
\]
La cubatura final de $S$ consumió 25,13 segundos de CPU. El cálculo a 2048 celdas de $P_T$ consumió
111,75 segundos de CPU en un hilo. Su semiancho es $0{,}0009121018263615$. Para el término de frontera
$J$ queda por integrar conjuntamente los coeficientes con signo sobre la posición y el estado de
disponibilidad del par.

\paragraph{Qué limita el encierro de $Q_2$.}
Para estudiar una de las filas con resolvente se eligió el nodo 10 del panel de 20 puntos. Su
peso de cuadratura dividido por $x^2$ vale aproximadamente $0{,}294$. Aun dedicando toda la meta
de $10^{-4}$ a ese nodo, el semiancho de $Q_2$ tendría que ser de unos $3{,}4\,10^{-4}$.
La integración conservó las dos paridades diagonales y el signo de los 17 aportes de traza tras
la cancelación. Se partió de un modelo certificado de $L_0$ con 256 paneles por unidad. Para
integrar se agruparon en 16 paneles por unidad y se usaron 512 centros. La integral completa del
continuo quedó en $[-49913{,}764054;50371{,}535812]$ [O], con un semiancho de unos $5{,}0\,10^4$.
El cálculo consumió 51,27 segundos de CPU.

¿Bastaría con desagrupar los paneles? En la traza $S_{13}$, las dos orientaciones
axiales entran con el mismo signo. Cada una contiene dos términos del bloque fila--columna de $L_0$,
con coeficiente uno. La menor anchura entre las $65536$ entradas de ese bloque fue $w=108349/131072$.
Las dos densidades axiales tienen masa uno. Tras multiplicar por el prefactor $1/4$, la evaluación
intervalar separada de esta traza deja un semiancho de al menos $w/2>0{,}41$, aun desagrupando todos
los paneles de origen. El mínimo de la tabla está certificado [O]. La cota inferior [D] vale para esa
tabla cuando se
evalúa $S_{13}$ por separado. Las identidades del núcleo físico pueden ligar más estrechamente
sus entradas.

Al reunir las cotas globales se obtiene $[-4{,}28\,10^{19};4{,}28\,10^{19}]$ [O]. El intervalo
incluye cero y deja abiertos tanto el signo como los decimales de $\cerr$. Para obtener un encierro
fino por la vía del método de rectángulos haría falta una cota demostrada y suficientemente pequeña
del resto espacial
en las clases de malla físicas. Así se podrían justificar las tolerancias espaciales de los nodos
a las que sigue condicionada la cota de integración en $x$.
